% for copilot. 
% figure file path mapping: 
% fig:ch2-lh2-schem: don't worry about this one, it's just a schematic of the system
% fig:coupled-ch2-lh2-t1-temp-dens maps to /Users/dmrraso/dv/victor-repos/hydrogen_fuel_tank/analysis/multi_tank_systems/coupled_ch2_lh2/output/plots/Coupled_CH2_LH2_Multi_Tank_System_(Feedforward)_evolution_density_temperature_tank1.png
% fig:coupled-ch2-lh2-t2-temp-dens maps to /Users/dmrraso/dv/victor-repos/hydrogen_fuel_tank/analysis/multi_tank_systems/coupled_ch2_lh2/output/plots/Coupled_CH2_LH2_Multi_Tank_System_(Feedforward)_evolution_density_temperature_tank2.png
% fig:coupled-ch2-lh2-t1-evolution maps to /Users/dmrraso/dv/victor-repos/hydrogen_fuel_tank/analysis/multi_tank_systems/coupled_ch2_lh2/output/plots/Coupled_CH2_LH2_Multi_Tank_System_(Feedforward)_evolution_evolution_tank1.png
% fig:coupled-ch2-lh2-t2-evolution maps to /Users/dmrraso/dv/victor-repos/hydrogen_fuel_tank/analysis/multi_tank_systems/coupled_ch2_lh2/output/plots/Coupled_CH2_LH2_Multi_Tank_System_(Feedforward)_evolution_evolution_tank2.png
% fig:coupled-ch2-lh2-t1-flows maps to /Users/dmrraso/dv/victor-repos/hydrogen_fuel_tank/analysis/multi_tank_systems/coupled_ch2_lh2/output/plots/Coupled_CH2_LH2_Multi_Tank_System_(Feedforward)_evolution_mass_flows_tank1.png
% fig:coupled-ch2-lh2-t2-flows maps to /Users/dmrraso/dv/victor-repos/hydrogen_fuel_tank/analysis/multi_tank_systems/coupled_ch2_lh2/output/plots/Coupled_CH2_LH2_Multi_Tank_System_(Feedforward)_evolution_mass_flows_tank2.png
% fig:ch2-lh2-req-press maps to /Users/dmrraso/dv/victor-repos/hydrogen_fuel_tank/analysis/multi_tank_systems/coupled_ch2_lh2/output/plots/Coupled_CH2_LH2_Multi_Tank_System_(Feedforward)_evolution_pressure_requirements.png




\subsubsection{A coupled \gls{ch2}-\gls{lh2} storage system} \

The coupled \gls{ch2}--\gls{lh2} system architecture is illustrated schematically in \autoref{fig:ch2-lh2-schem}, showing the high-pressure \gls{ch2} tank connected to the low-pressure \gls{lh2} tank through a \gls{pid}-controlled variable flow valve that provides mission-adaptive pressurisation.

\begin{figure}[H]
    \centering
     \includegraphics[width=0.8\linewidth]{\figpath ch2_lh2_schem.jpg}
    \caption{Bi-tank architecture for a coupled \gls{ch2}-\gls{lh2} storage design}
    \label{fig:ch2-lh2-schem}
\end{figure}

The key properties defining each tank in the coupled \gls{ch2}--\gls{lh2} system are summarised in \autoref{tab:ch2-lh2-tanks}. These inputs specify the initial conditions, materials, and heat-transfer parameters used for both tanks. The \gls{ch2} reservoir (tank~1) functions as the high-pressure pressurisation source, while the \gls{lh2} tank (tank~2) is the low-pressure fuel supplier that interfaces directly with the ATR72 mission discharge requirements.

\begin{table}[H]
\centering
\caption{Input parameters for the coupled \gls{ch2}--\gls{lh2} system. Tank~1 serves as the high-pressure pressurisation source, and tank~2 supplies mission fuel.}
\label{tab:ch2-lh2-tanks}
\begin{tabular}{@{}lcc@{}}
\toprule
\textbf{Parameter} & \textbf{Tank 1 (\gls{ch2})} & \textbf{Tank 2 (\gls{lh2})} \\ \midrule
Storage type & \gls{ch2} & \gls{lh2} \\
$P_{\text{init}}$ [bar] & 1000 & 10 \\
$T_{\text{init}}$ [K] & 288.1 & 23.5 \\
$\rho_{\text{init}}$ [kg/m$^3$] & 60.0 & 68.0 \\
$P_{\text{vent}}$ [bar] & 1100 & 15 \\
$P_{\text{min}}$ [bar] & 1 & 2 \\
Tank volume [m$^3$] & 2.145 & 3.230 \\
Tank radius [m] & 0.80 & 0.917 \\
Wall material [-] & Carbon-epoxy & Carbon-epoxy \\
Liner material [-] & Aluminium 6061-T6 & Aluminium 6061-T6 \\
Liner thickness [mm] & 1.0 & 1.0 \\
Insulation \gls{htc} [W/m$^2$K] & 50 & 0.025 \\
Sizing method [-] & User-defined radius & Mission-based \\
Discharge rate [kg/s] & Determined by coupling logic & ATR72 mission profile \\ \bottomrule
\end{tabular}
\end{table}

The operational logic of the mission-adaptive pressurisation system and its governing control parameters are listed in \autoref{tab:ch2-lh2-coupling}. Unlike simple pressure-triggered valves, this configuration employs a \gls{pid}-controlled variable flow valve that responds continuously to mission demands. The \gls{pid} controller modulates the pressurisation mass flow rate in real time based on feedback from tank~2 pressure measurements. The parameters listed are demonstrative rather than optimised, serving to illustrate how the coupling rule is defined and executed within the framework.

\begin{table}[H]
\centering
\caption{\gls{pid}-controlled mission-adaptive pressurisation configuration for the coupled \gls{ch2}--\gls{lh2} storage system.}
\label{tab:ch2-lh2-coupling}
\begin{tabular}{@{}lc@{}}
\toprule
\textbf{Parameter} & \textbf{Value / Description} \\ \midrule
Coupling type & Mission-adaptive pressurisation \\
Control logic & \gls{pid}-controlled variable flow valve \\
Pressure margin [bar] & 0.5 (above minimum required) \\
Minimum source pressure [bar] & 15.0 \\
Maximum pressurisation flow rate [kg/s] & 0.11 \\ \midrule
\multicolumn{2}{c}{\textbf{\gls{pid} controller parameters}} \\ \midrule
Proportional gain, $K_p$ [-] & 0.8 \\
Integral gain, $K_i$ [-] & 0.08 \\
Derivative gain, $K_d$ [-] & 0.1 \\
Deadband control [bar] & 0.05 \\ \midrule
\multicolumn{2}{c}{\textbf{Discharge physics}} \\ \midrule
Discharge piping diameter [mm] & 5.0 \\
Discharge piping length [m] & 1.0 \\
Discharge coefficient [-] & 0.6 \\
Pipe surface roughness [m] & $1.5 \times 10^{-6}$ (stainless steel) \\
Loss coefficient [-] & 1.0 (minimal minor losses) \\ \midrule
\multicolumn{2}{c}{\textbf{Coupling physics configuration}} \\ \midrule
Fluid properties & CoolProp (H$_2$) \\
Choked flow & Enabled for demonstration \\
Reynolds transition [-] & 2300 (laminar to turbulent) \\
Orifice diameter [mm] & 8.0 \\
Friction correlation [-] & Blasius \cite{blasius} \\
Maximum velocity limit [m/s] & 2000 \\ \bottomrule
\end{tabular}
\end{table}

\begin{figure}[H]
    \centering
    \begin{subfigure}[b]{0.45\textwidth}
         \includegraphics[width=\linewidth]{\figpath coupled_ch2_lh2_t1_temp_dens.jpg}
        \caption{Density–temperature plot for tank~1 (\gls{ch2}) in a coupled \gls{ch2}-\gls{lh2} storage system providing fuel for the ATR72 mission.}
        \label{fig:coupled-ch2-lh2-t1-temp-dens}
    \end{subfigure}
    \hfill
    \begin{subfigure}[b]{0.45\textwidth}
        \includegraphics[width=\linewidth]{\figpath coupled_ch2_lh2_t2_temp_dens.jpg}
        \caption{Density–temperature plot for tank~2 (\gls{lh2}) in a coupled \gls{ch2}-\gls{lh2} storage system providing fuel for the ATR72 mission.}
        \label{fig:coupled-ch2-lh2-t2-temp-dens}
    \end{subfigure}
    \caption{}
\end{figure}

At the start of the mission, the valve opens immediately due to the high flow demand during take-off. The resulting pressure deficit in tank~2 triggers an immediate transfer of high-pressure, warm hydrogen from tank~1, as seen in the steep initial gradients of the trajectories in \autoref{fig:coupled-ch2-lh2-t1-temp-dens} and \autoref{fig:coupled-ch2-lh2-t2-temp-dens}. The injection of warm hydrogen into the cryogenic tank rapidly pushes the \gls{lh2} tank into a two-phase regime—a condition that would be undesirable in a practical implementation but is inevitable in this idealised demonstration owing to the sensitivity of liquid hydrogen to boil-off. Persistent oscillations in the plotted variables throughout the mission reflect the action of the \gls{pid} controller as it continuously adjusts the flow rate to maintain the desired pressure range. These fluctuations are a natural response of the closed-loop system to transient mass-flow perturbations.

\begin{figure}[H]
    \centering
    \begin{subfigure}[b]{0.45\textwidth}
         \includegraphics[width=\linewidth]{\figpath coupled_ch2_lh2_t1_evolution.jpg}
        \caption{State evolution for tank~1 (\gls{ch2}) in a coupled \gls{ch2}-\gls{lh2} storage system providing fuel for the ATR72 mission.}
        \label{fig:coupled-ch2-lh2-t1-evolution}
    \end{subfigure}
    \hfill
    \begin{subfigure}[b]{0.45\textwidth}
        \includegraphics[width=\linewidth]{\figpath coupled_ch2_lh2_t2_evolution.jpg}
        \caption{State evolution for tank~2 (\gls{lh2}) in a coupled \gls{ch2}-\gls{lh2} storage system providing fuel for the ATR72 mission.}
        \label{fig:coupled-ch2-lh2-t2-evolution}
    \end{subfigure}
    \caption{}
\end{figure}

The pressure, temperature, and density evolutions shown in \autoref{fig:coupled-ch2-lh2-t1-evolution} and \autoref{fig:coupled-ch2-lh2-t2-evolution} display complementary behaviour between the two tanks. The pressure reduction in tank~1 coincides with a corresponding increase in tank~2, validating the conservation of mass. The initial rapid pressure drop in the \gls{lh2} tank corresponds to the steep fuel extraction rate during take-off; as the controller activates, the pressurisation flow compensates for this deficit and restores the pressure within the operational limits for a typical \gls{lh2} storage system. Both tanks retain a notable quantity of hydrogen at mission end, suggesting that smaller tank could be used to complete the same mission.

\begin{figure}[H]
    \centering
    \begin{subfigure}[b]{0.45\textwidth}
         \includegraphics[width=\linewidth]{\figpath coupled_ch2_lh2_t1_flow_rates.jpg}
        \caption{Flow rates for tank~1 (\gls{ch2}) in a coupled \gls{ch2}-\gls{lh2} storage system providing fuel for the ATR72 mission.}
        \label{fig:coupled-ch2-lh2-t1-flows}
    \end{subfigure}
    \hfill
    \begin{subfigure}[b]{0.45\textwidth}
        \includegraphics[width=\linewidth]{\figpath coupled_ch2_lh2_t2_flow_rates.jpg}
        \caption{Flow rates for tank~2 (\gls{lh2}) in a coupled \gls{ch2}-\gls{lh2} storage system providing fuel for the ATR72 mission.}
        \label{fig:coupled-ch2-lh2-t2-flows}
    \end{subfigure}
    \caption{}
\end{figure}

The coupling flow-rate profiles shown in \autoref{fig:coupled-ch2-lh2-t1-flows} and \autoref{fig:coupled-ch2-lh2-t2-flows} further confirm the transient nature of the pressurisation process. A high initial coupling flow is followed by a gradual reduction in flow magnitude as the mission progresses. The continuing boil-off and mixing with warm hydrogen raises the mean temperature in tank~2 and consequently, the pressure compensation required from tank~1 decreases accordingly, leading to a steady decline in transfer rate. The plateau observed in the first flow spike indicates that the flow limit was reached, implying that the valve was operating at its maximum allowable throughput. A more optimised valve design or a tuned control law could yield smoother and more efficient flow profiles. It should also be noted that the model disables choked-flow physics for simplicity; including such effects would slightly alter the transient profiles.

\begin{figure}[H]
    \centering
     \includegraphics[width=0.8\linewidth]{\figpath coupled_ch2_lh2_req_press.jpg}
    \caption{Comparison of actual pressure and activation pressure in tank~2 for the control law in the \gls{ch2}-\gls{lh2} system.}
    \label{fig:ch2-lh2-req-press}
\end{figure}

\autoref{fig:ch2-lh2-req-press} compares the instantaneous pressure in tank~2 with the corresponding target activation pressure defined by the control law. The \gls{pid} controller strives to ensure that the actual pressure (solid line) tracks the desired trajectory (grey markers). Early in the mission, the controller exhibits some overshoot and undershoot due to the abrupt fuel demand associated with take-off, which challenges the control authority of the system. At the very start of the mission, the controller cannot compensate for the minimum required pressure in tank 2 quickly enough, but then quickly recovers. Throughout the mission, however, the actual pressure in tank 2 never falls below the absolute minimum of 2 bar. The gains used here were intentionally left unoptimised; appropriate tuning would achieve a smoother and more stable response. The selected 1\,bar pressure margin between the minimum and target activation pressures is arbitrary and serves purely to demonstrate the controller’s feedback operation.

Overall, the coupled \gls{ch2}–\gls{lh2} configuration successfully demonstrates how active control logic can be incorporated into the sizing tool framework. The use of a \gls{pid}-based pressurisation strategy shows that complex, dynamic coupling rules can be embedded within the same modular architecture used for simpler systems. Although this configuration is not suitable as a practical design—owing to the large size and high pressure required of the \gls{ch2} tank and the two-phase behaviour induced in the \gls{lh2} tank—it nonetheless highlights the adaptability of the modelling framework. 